\section{Step 20 Advance tile origins in steps of four}\label{transpose-s20}

After completing a tile, advance its first column by four. If another tile column begins inside the matrix, calculate its clipped ends and process it. At the end of a tile row, advance the tile's first row by four and reset the first column through \texttt{next\_tile\_row}. Cell counters advance by one; tile origins advance by four.

\subsection{Write or inspect}
Append the outer tile loops. Run to the return boundary with three by five, then five by three and five by nine.

\Needspace{11\baselineskip}
\MARSCodeLabel{MIPS assembly}
\begin{minted}[samepage=true]{mars-idea}
    addiu $t1, $t1, 4
    sltu $t9, $t1, $a3
    bne $t9, $zero, next_tile
    nop
    addiu $t0, $t0, 4
    sltu $t9, $t0, $a2
    bne $t9, $zero, next_tile_row
    nop
\end{minted}

\subsection{Expected observation}
The default second tile copies four, nine, and fourteen. Both functions now produce the same five-by-three output while following different access orders.

\Needspace{11\baselineskip}
\subsection{Check your understanding}
Why must the tile column origin reset when the tile row changes?

\begin{itemize}
  \item \textbf{A} To clear the output cache.
  \item \textbf{B} To increase the return count.
  \item \textbf{C} To cover the leftmost tile in the next group of rows.
\end{itemize}
